% Source presentations for register row N-c-026.

\subsection{Source-exact complexity boundaries for SAT, QBF, and modal logic}
\label{sec:p09-presentation-n-c-026}

\subsubsection{Cook's historical reduction and the modern SAT theorem}

\paragraph{The reduction Cook actually defines.}
A \emph{query machine} is a deterministic multitape Turing machine with a
query tape and three distinguished states: query, yes, and no.  Relative to a
language \(T\), entering the query state with \(u\) on the query tape moves in
one step to yes when \(u\in T\) and to no otherwise.  Cook calls a language
\(S\) \emph{\(P\)-reducible} to \(T\) when such a machine decides \(S\) in
time polynomial in its input length.  The term therefore stands for an
oracle computation that may issue several adaptive queries.  It is not the
single-output many-one map fixed in Part~IV, and this distinction is used when
the historical theorem is related to its modern descendant.

Cook's Theorem~1 says that if a nondeterministic Turing machine accepts \(S\)
in polynomial time, then \(S\) is \(P\)-reducible in this query-machine sense
to the disjunctive-normal-form tautologies.  For input \(w\), his proof first
constructs a conjunctive-normal-form formula \(A(w)\) satisfying
\[
 w\in S
 \quad\Longleftrightarrow\quad
 A(w)\text{ is satisfiable}.
\]
Consequently \(\neg A(w)\) is already in DNF and
\[
 \neg A(w)\text{ is a tautology}
 \quad\Longleftrightarrow\quad
 w\notin S.
\]
The query machine asks the one question \(\neg A(w)\) and accepts when the
oracle answers no.  This proves Cook's stated query reduction.  The displayed
map \(w\mapsto A(w)\), viewed separately, is the bridge to the later many-one
SAT theorem \cite[pp.~151--156]{Cook1971}.

\paragraph{A complete modern reconstruction of the tableau.}
Let \(L\in\mathsf{NP}\), and fix a nondeterministic machine deciding \(L\) in
polynomial time.  A polynomial-length string \(y\) can record its successive
nondeterministic choices.  A fixed deterministic machine \(V\) then checks the
computation encoded by \((x,y)\) within a polynomial time bound \(T=T(|x|)\),
and
\[
 x\in L
 \quad\Longleftrightarrow\quad
 \exists y\; V(x,y)\text{ accepts within }T\text{ steps}.
\]
Invalid choice codes are made rejecting, and the accepting state is made
absorbing, so it is enough to inspect the configuration at exactly time \(T\).
This deterministic verifier presentation is used only to make every local
tableau successor uniquely determined once the free witness bits have been
chosen.  A standard one-tape simulation is fixed once and for all; its
polynomial overhead is absorbed into \(T\).

Let \(\Sigma\), including fixed nonoverwritable endmarkers, and \(Q\) be the
tape alphabet and state set of \(V\).
A tableau cell carries a label in the fixed finite alphabet
\[
 \Lambda=\Sigma\cup(Q\times\Sigma).
\]
The label \(\sigma\) says that the cell contains \(\sigma\) and is not scanned;
the label \((q,\sigma)\) additionally says that the head scans that cell while
the machine is in state \(q\).  This packaging is an eliminable abbreviation
for Cook's separate tape-symbol, state, and scanned-square proposition
families; it changes no computation condition and is used to make locality
visible.

Choose a width \(W\) large enough to hold the initial encoding of \(x\) and
\(y\), both endmarkers, and every cell reachable in \(T\) steps.  Thus
\(W=O(T+|x|+|y|)\).  For
\(0\leq t\leq T\), \(1\leq s\leq W\), and
\(\lambda\in\Lambda\), introduce a propositional variable
\(X_{t,s,\lambda}\), intended to say that cell \(s\) has label \(\lambda\) at
time \(t\).

For a finite list of formulas \(Z_1,\ldots,Z_k\), use the metalinguistic
abbreviation
\[
 \operatorname{One}(Z_1,\ldots,Z_k)
   = (Z_1\lor\cdots\lor Z_k)
     \land
     \bigwedge_{1\leq i<j\leq k}(\neg Z_i\lor\neg Z_j).
\]
It is a CNF formula, not a new connective: the first clause supplies at least
one true member and the binary clauses forbid two.  It is needed twice below.
For every \((t,s)\), include
\[
 \operatorname{One}\bigl((X_{t,s,\lambda})_{\lambda\in\Lambda}\bigr),
\]
and, for every row \(t\), require exactly one variable whose label lies in
\(Q\times\Sigma\).  Every satisfying assignment therefore decodes to one tape
symbol at every position and exactly one state/head position in every row.

The remaining clauses have three roles.

\begin{enumerate}
\item The \emph{initial clauses} fix the endmarkers, the literal bits of \(x\),
the initial state and head position, and the blank cells.  Each witness position
is restricted to the two labels for \(0\) and \(1\) but is otherwise free, so a
satisfying assignment chooses exactly one string \(y\).

\item The \emph{transition clauses} express the fixed finite transition table
of \(V\).  With fixed boundary labels, the next label at position \(s\) depends
only on the three labels at positions \(s-1,s,s+1\) in the preceding row.  Let
\(\Delta\subseteq\Lambda^3\times\Lambda\) be exactly the finite set of allowed
quadruples.  For every \(t<T\), every interior \(s\), and every
\((\alpha,\beta,\gamma,\delta)\notin\Delta\), include
\[
 \neg X_{t,s-1,\alpha}\lor
 \neg X_{t,s,\beta}\lor
 \neg X_{t,s+1,\gamma}\lor
 \neg X_{t+1,s,\delta}.
\]
The formula forbids precisely a next-cell label incompatible with the local
window.  Boundary positions use the analogous fixed-endmarker clauses.

\item The \emph{acceptance clause} is
\[
 \bigvee_{1\leq s\leq W,\,\sigma\in\Sigma}
       X_{T,s,(q_{\mathrm{acc}},\sigma)}.
\]
Because the accepting state is absorbing, this says exactly that the encoded
run accepts within the original time bound.
\end{enumerate}

Let \(F_{V,x}\) be the conjunction of these clauses.  If \(x\in L\), choose an
accepting witness \(y\) and give each \(X_{t,s,\lambda}\) its actual tableau
value.  Every exactly-one, initial, transition, and acceptance clause is true,
so \(F_{V,x}\) is satisfiable.  Conversely, a satisfying assignment first
selects one label for every cell and one head/state label per row.  The initial
clauses produce the genuine initial configuration of \(V(x,y)\) for the chosen
\(y\).  Induction on \(t\), using the forbidden local quadruples, shows that
each later row is its genuine successor.  The final clause therefore supplies
an accepting computation.  Hence
\[
 x\in L
 \quad\Longleftrightarrow\quad
 F_{V,x}\in\mathsf{SAT}.
\]

The label alphabet is fixed with \(V\).  There are \(O(TW)\) tableau variables;
the cell and transition clauses have polynomial total size, and even the naive
pairwise exactly-one condition for each row has only \(O(TW^2)\) clauses.
Nested counters over \(t,s\) and the fixed label alphabet emit the formula in
polynomial time.  Thus \(x\mapsto F_{V,x}\) is a polynomial-time many-one
reduction.

Finally, \(\mathsf{SAT}\in\mathsf{NP}\): an assignment has at most one bit per
variable occurring in the encoded formula, and its truth value is checked by
one traversal of the formula.  The reduction and this membership argument
give the exact modern statement in Theorem~\ref{thm:p04-cook-levin}; the
written proof under the machine conventions of Part~4, including the passage
from nondeterministic machines to certificates, is
Annex~\ref{annex:p02-w02b-p02-015-cook-levin}.
Cook's 1971 paper controls the historical query-reduction statement; Levin's
independent 1973 work supplies the other historical source for the compound
name \cite{Cook1971,Levin1973,AroraBarak2009}.

\subsubsection{Stockmeyer--Meyer's quantified language and the complete proof}

\paragraph{The source language and the scalar-variable bridge.}
Stockmeyer and Meyer write
\(X_i=(x_{i,1},x_{i,2},\ldots)\) for a finite block of Boolean variables.  For
\(k\geq1\), their language \(B_k\) consists of encoded Boolean expressions
\(A(X_1,\ldots,X_k)\) made true by the alternating prefix
\[
 \exists X_1\,\forall X_2\,\exists X_3\cdots Q_kX_k,
 \qquad
 Q_k=
 \begin{cases}
  \exists,&k\text{ odd},\\
  \forall,&k\text{ even}.
 \end{cases}
\]
The notation \(\exists X_i\) abbreviates one Boolean quantifier for each scalar
variable in the encoded block; it is not a quantifier over an infinite object.
The union
\[
 B_\omega=\bigcup_{k\geq1}B_k
\]
allows the finite alternation depth to grow with the input.  It is one language
of true encoded formulas, not the polynomial hierarchy as a class of languages.
The source's Theorem~4.3 states that \(B_\omega\) is complete for polynomial
space under logspace many-one reductions
\cite[thm.~4.3]{StockmeyerMeyer1973}.

The bridge to Definition~\ref{def:p04-qbf-languages} is syntactic.  Expand each
block quantifier into its scalar Boolean quantifiers and retain the same Boolean
matrix.  Conversely, rename and collect every maximal run of like scalar
quantifiers into one block \(X_i\); if a TQBF begins universally, prepend one
unused existential variable so that the block sequence begins existentially.
If the chosen scalar encoding admits quantifiers inside Boolean subformulas,
first rename bound variables apart, push negations through quantifiers, and pull
each quantifier across a neighboring subformula in which its variable is not
free.  These equivalences produce prenex form without copying the matrix.
For already prenex inputs the two block/scalar maps have linear output length;
the optional prenex normalization remains polynomial in length.  In either
case a transducer needs only an input-position counter, a block counter, and a
polarity bit, so it uses logarithmic work space.  No compilation of an unstated
polynomial-time predicate is required: \(B_k\) already contains the Boolean
expression as part of its input.

\paragraph{Membership of TQBF in polynomial space.}
Evaluate a closed QBF recursively.  At \(\exists p\), evaluate the body with
\(p=0\) and then, only if necessary, with \(p=1\); at \(\forall p\), require
both evaluations.  The two branches are explored sequentially.  Store the
current syntax position, the values of variables on the active quantifier path,
and one return bit per recursive level.  Each object has size polynomial in the
input formula, so the procedure may take exponential time but uses polynomial
space.  Thus \(\mathsf{TQBF}\in\mathsf{PSPACE}\).

\paragraph{Polynomial-space hardness without recursive duplication.}
Let a deterministic machine \(M\) decide a language in polynomial space.
Modify it so that acceptance first reaches one canonical accepting
configuration.  On input \(x\), encode every configuration by an \(m\)-bit
block, where \(m\geq1\) is polynomial in \(|x|\).  There are Boolean formulas
\(\operatorname{Eq}(u,v)\) and
\(\operatorname{Step}_{M,x}(u,v)\) of polynomial size saying, respectively,
that two blocks agree bitwise and that \(v\) is the unique legal successor of a
valid \(u\).  Put
\[
 R_0(u,v)=\operatorname{Eq}(u,v)\lor
          \operatorname{Step}_{M,x}(u,v).
\]
Here and below \(R_i\) is a metalinguistic family of ordinary QBFs with two free
\(m\)-bit parameter blocks.  It adds no reachability connective to the
language.

If the naive recursion conjoined \(R_i(u,z)\) and \(R_i(z,v)\), its text would
double at every level.  Instead, for fresh \(m\)-bit blocks \(z,a,c\) and a
fresh selector bit \(b\), define
\[
\begin{aligned}
 \operatorname{Sel}_b(u,z,v;a,c)
  ={}&\bigl(\neg b\to
       (\operatorname{Eq}(a,u)\land\operatorname{Eq}(c,z))\bigr)\\
    &\land\bigl(b\to
       (\operatorname{Eq}(a,z)\land\operatorname{Eq}(c,v))\bigr),\\
 R_{i+1}(u,v)
  ={}&\exists z\,\forall b\,\exists a\,\exists c\,
       \bigl(\operatorname{Sel}_b(u,z,v;a,c)\land R_i(a,c)\bigr).
\end{aligned}
\]
Every displayed block quantifier expands into \(m\) scalar Boolean quantifiers.
The existential blocks \(a,c\) follow \(b\), so they may receive the first pair
\((u,z)\) when \(b=0\) and the second pair \((z,v)\) when \(b=1\).  There is
only one textual occurrence of \(R_i\), and hence
\[
 |R_{i+1}|=|R_i|+O(m).
\]

Induction on \(i\) proves, for all \(m\)-bit blocks \(u,v\),
\[
 R_i(u,v)\text{ is true}
 \quad\Longleftrightarrow\quad
 v\text{ is reachable from }u\text{ by at most }2^i
 \text{ equality-or-machine steps}.
\]
The base is the definition of \(R_0\).  At the successor step, a true formula
supplies one midpoint and the universal selector establishes both half paths;
concatenating them gives a path of length at most \(2^{i+1}\).  Conversely,
split any such path after at most \(2^i\) steps.  The two pieces satisfy the
induction hypothesis, and the selector routes their endpoints to the single
recursive occurrence.  Equality steps merely pad a shorter path.

There are at most \(2^m\) encoded configurations.  A deterministic accepting
run has no repeated preaccepting configuration, because a repetition would
begin a loop and prevent later acceptance.  Therefore
\[
 \Phi_{M,x}=R_m(c_{\mathrm{init}}(x),c_{\mathrm{acc}})
\]
is true exactly when \(M\) accepts \(x\).  Its length is polynomial: the base
formula is polynomial and each of the \(m\) recursion levels adds only
\(O(m)\) symbols.

The map is also logspace-constructible, as in the historical source strength.
A transducer stores a level counter and the indices of the currently named
blocks, emits the \(m\) nested prefixes from outside inward, streams the fixed
local clauses for \(\operatorname{Step}_{M,x}\) with logarithmic counters, and
then emits the closing suffixes in the reverse level order.  All names and
counters require \(O(\log m)=O(\log|x|)\) bits because the machine and its
polynomial space bound are fixed.  This proves PSPACE-hardness under logspace
many-one reductions, and therefore proves the weaker polynomial-time
many-one statement of Theorem~\ref{thm:p04-tqbf-pspace}.  The full formal
argument, including validity of blocks that do not encode configurations, is
Annex~\ref{annex:p9.tqbf.1}.

\paragraph{Papadimitriou's distinct QSAT presentation.}
Papadimitriou's printed p.~455 uses \emph{QSAT} for the language of true closed
formulas
\[
 \exists x_1\,\forall x_2\,\exists x_3\cdots Q_n x_n\,
 \varphi(x_1,\ldots,x_n),
\]
where the scalar Boolean quantifiers alternate strictly and \(\varphi\) is in
conjunctive normal form.  This is a decision language, not an abbreviation for
the unrestricted input syntax of Definition~\ref{def:p04-qbf-languages}; it
also differs from Stockmeyer--Meyer's block-variable \(B_\omega\).  The
distinction is needed to state the book's theorem at its own strength.

The source contains an internal conflict that must not be normalized away.
The display, the parity clause for \(Q_n\), and the proof's alternating
OR/AND evaluation all begin existentially, but the prose immediately before
the display describes both truth values for \(x_1\), which would begin
universally.  The inspected errata give no correction for p.~455.  The
display-and-proof-controlled reading above is therefore the coherent one, but
the contradictory prose remains part of the source record.

Theorem~19.1, printed pp.~456--458, states that this QSAT language is
\(\mathsf{PSPACE}\)-complete.  Under the book's standing definition, the
reduction is a deterministic logspace-computable membership-preserving
function---a logspace many-one reduction in current terminology.  Membership
evaluates the implicit alternating assignment tree depth first.  For hardness,
the proof converts polynomial-space configuration reachability into a
single-copy formula family.  Reusing the bitwise equality and legal-successor
formulas defined above, put
\[
 \psi_0(A,B)=\operatorname{Eq}(A,B)\lor
              \operatorname{Step}_{M,x}(A,B).
\]
For fresh blocks \(Z,X,Y\) at each level, the recurrence has the eliminable
form
\[
 \psi_{i+1}(A,B)=
 \exists Z\,\forall X\,\forall Y\,
 \left(
   \bigl((X=A\land Y=Z)\lor(X=Z\land Y=B)\bigr)
   \to \psi_i(X,Y)
 \right),
\]
where each capital letter denotes a finite Boolean block and each displayed
block equality expands bitwise.  The \(\psi_i\) are a metalinguistic family of
ordinary quantified formulas, not a new reachability connective.  This
recurrence is recorded separately because its universally quantified endpoint
blocks are not the selector-bit construction above; it is used here only to
preserve Papadimitriou's own presentation.  After prenexing, the proof handles
the polynomial DNF matrix through the complement of the displayed CNF language
and \(\mathsf{PSPACE}=\mathsf{coPSPACE}\), rather than silently asserting that
the prenexed matrix is already a polynomial CNF
\cite[pp.~455--458]{Papadimitriou1994}.

Google Books primary-text snippets independently confirm the p.~455 definition
and conflict, the p.~456 theorem and proof opening, and the p.~458 logspace and
complement conclusion.  Its volume remains nonviewable, however, and no
readable publisher or named-library body was found for the continuous p.~457
recurrence or the p.~487 historical note.  The latter says that QSAT
PSPACE-completeness was established in the Stockmeyer--Meyer paper; it does not
literally call that paper the ``first proof.''  Those two page bodies therefore
remain provenance-qualified even though the theorem statement, locator,
intended syntax, reduction strength, and earlier first-party proof lineage are
independently secured.

\subsubsection{Ladner's polarity-correct modal table}

Ladner's input is a modal propositional formula together with a fixed choice of
system, not a supplied finite Kripke model.  Under the Kripke semantics used in
the source, \(\mathsf K\) has arbitrary frames, \(\mathsf T\) reflexive frames,
\(\mathsf S4\) reflexive-transitive frames, and \(\mathsf S5\) equivalence
frames.  For each of these systems, Kripke completeness identifies provability
with validity.  Negation gives the distinct dualities
\[
 \begin{aligned}
  A\text{ is not valid}
    &\quad\Longleftrightarrow\quad \neg A\text{ is satisfiable},\\
  A\text{ is not provable}
    &\quad\Longleftrightarrow\quad \neg A\text{ is satisfiable}.
 \end{aligned}
\]
The second equivalence uses completeness.  It does not identify nonprovability
of \(A\) with satisfiability of the same formula \(A\); the linear negation map
is part of the reduction.  With that polarity exposed, the exact rows are:

\begin{table}[htbp]
\centering\scriptsize
\begin{tabularx}{\textwidth}{@{}l l X X@{}}
\toprule
System & Frame class & Decision problem & Exact class and reduction \\
\midrule
\(\mathsf K\) & arbitrary & provability/validity &
\(\mathsf{PSPACE}\)-complete under logspace many-one reductions \\
\(\mathsf K\) & arbitrary & satisfiability; equivalently nonprovability after
negating the input & \(\mathsf{PSPACE}\)-complete under the same reductions \\
\(\mathsf T\) & reflexive & provability/validity &
\(\mathsf{PSPACE}\)-complete under logspace many-one reductions \\
\(\mathsf T\) & reflexive & satisfiability; equivalently nonprovability after
negating the input & \(\mathsf{PSPACE}\)-complete under the same reductions \\
\(\mathsf S4\) & reflexive and transitive & provability/validity &
\(\mathsf{PSPACE}\)-complete under logspace many-one reductions \\
\(\mathsf S4\) & reflexive and transitive & satisfiability; equivalently
nonprovability after negating the input & \(\mathsf{PSPACE}\)-complete under
the same reductions \\
\(\mathsf S5\) & equivalence & satisfiability; equivalently nonprovability
after negating the input & \(\mathsf{NP}\)-complete under logspace many-one
reductions \\
\(\mathsf S5\) & equivalence & provability/validity &
\(\mathsf{coNP}\)-complete under logspace many-one reductions \\
\bottomrule
\end{tabularx}
\caption{Ladner's modal complexity results with decision polarity exposed.}
\label{tab:p09-ladner-polarity}
\end{table}

\paragraph{The common lower bound from quantified Boolean truth.}
Ladner calls a modal system \(S\) \emph{between \(\mathsf K\) and
\(\mathsf S4\)} when
\[
 \mathsf K\text{-PROVABLE}
 \subseteq S\text{-PROVABLE}
 \subseteq \mathsf S4\text{-PROVABLE}.
\]
This is an inclusion condition on theorem sets, not a new frame class.
Theorem~3.1 proves that \(B_\omega\) is logspace many-one reducible to
\(S\)-provability for every such \(S\).  The explicit completeness conclusions
later combine this lower bound with upper algorithms only for
\(\mathsf K,\mathsf T,\mathsf S4\).

The reduction first turns a closed QBF
\[
 A=Q_1X_1\cdots Q_mX_m\,A'(X_1,\ldots,X_m)
\]
into a modal conjunction \(B\).  Auxiliary propositions mark quantifier levels
and record the values already chosen.  At a universal level the formula forces
accessible successors for both truth values; at an existential level one
appropriate successor suffices.  Box formulas propagate every earlier choice,
and the final-level obligation requires \(A'\).

The proof invariant at a world representing level \(i\) is that the recorded
values of \(X_1,\ldots,X_i\) make the remaining quantified suffix
\[
 Q_{i+1}X_{i+1}\cdots Q_mX_m\,A'
\]
true.  If \(A\) is true, its existential choices and universal branches build a
reflexive-transitive model satisfying \(B\).  Conversely, from any arbitrary
\(\mathsf K\)-model satisfying \(B\), the forced successors recover choices
that establish the invariant backwards and hence show that \(A\) is true.
The modal gadgets realizing the level markers and propagation occupy the long
middle of Ladner's Theorem~3.1 and remain cited rather than silently replaced
by this invariant.

At the source's exact polarity, the direct construction establishes
\[
 A\in B_\omega
 \quad\Longleftrightarrow\quad
 \neg B\notin S\text{-PROVABLE}.
\]
Ladner's preceding Fact~2.7 supplies a logspace self-complementing reduction
for \(B_\omega\); composing it with this construction yields the stated
reduction to \(S\)-PROVABLE rather than merely to its complement.  This is why
the lower bound cannot be summarized safely by dropping either the negation or
the decision polarity \cite[thm.~3.1]{Ladner1977}.

\paragraph{Why \(\mathsf K,\mathsf T,\mathsf S4\) use only polynomial space.}
Ladner's Section~5 reformulates Kripke satisfiability procedures as recursive
world-construction algorithms.  A call retains finite sets of subformulas that
must be true or false at the current world, checks their Boolean consistency,
and recursively constructs only the successor demanded by the next modal
obligation.  A satisfying model may have exponentially many worlds, but the
algorithm explores one branch at a time and reuses the space of a completed
branch.

The frame conditions determine what must be carried into the recursive call.
The \(\mathsf K\) procedure carries the direct box obligations;
\(\mathsf T\) also enforces them at the current reflexive world; and
\(\mathsf S4\) propagates them down the transitive branch.  Subformulas and the
propagated finite sets bound the recursion depth polynomially.  The source's
explicit bounds are
\[
 \mathsf K\text{-PROVABLE}\in\operatorname{SPACE}(n^2),\qquad
 \mathsf T\text{-PROVABLE}\in\operatorname{SPACE}(n^3),\qquad
 \mathsf S4\text{-PROVABLE}\in\operatorname{SPACE}(n^4).
\]
These are sufficient upper bounds, not optimality claims.  Together with
Theorem~3.1 they give the three PSPACE-completeness rows in
Table~\ref{tab:p09-ladner-polarity}; closure of PSPACE under complement and the
linear negation map give the corresponding satisfiability rows
\cite[thm.~5.1 and cor.~5.2--5.3]{Ladner1977}.

\paragraph{Why \(\mathsf S5\) has a smaller certificate.}
Suppose \(A\) is \(\mathsf S5\)-satisfiable and contains \(m\) modal
connectives.  Restrict a satisfying model to the equivalence class of its
distinguished world, so accessibility is universal on that class.  Keep the
distinguished world and, for each relevant false box occurrence \(\Box C\),
one accessible witness at which \(C\) is false.  True boxes remain true after
worlds are removed, and every retained false box keeps its witness.  Induction
over the selected subformula occurrences therefore preserves their truth
values.  At most one new world is needed per modal occurrence, so the retained
model has at most \(m+1\) worlds.  This is Ladner's Lemma~6.1, not a finite-model
claim for the other three systems.

A nondeterministic machine can guess the at-most-\(m+1\) worlds and the truth
values of the proposition letters occurring in \(A\), then evaluate every
subformula at every world in polynomial time.  Hence \(\mathsf S5\)
satisfiability is in NP.  Propositional SAT embeds as the fragment with no modal
operator, giving NP-hardness under logspace reductions.  Thus satisfiability,
and equivalently nonprovability after negation, is NP-complete; validity and
provability are its coNP-complete complements
\cite[lem.~6.1 and thm.~6.2]{Ladner1977}.  Saying only that “\(\mathsf S5\) is
NP-complete” would leave out the decision problem and would be false for the
validity/provability row unless \(\mathsf{NP}=\mathsf{coNP}\).

Modal model checking on a supplied finite Kripke structure is a different input
problem and does not inherit this provability table.  Likewise,
Theorem~\ref{thm:fol-validity-undecidable} says unrestricted first-order
validity is undecidable, while
Corollary~\ref{thm:fol-validity-enumerable} says it is recursively enumerable.
Those first-order conclusions neither give a finite-model result nor one of the
modal complexity classes above.
